\documentclass[11pt]{article}

% ===================== PACKAGES =====================
\usepackage[a4paper,margin=1in]{geometry}
\usepackage{amsmath,amssymb,amsthm}
\usepackage{enumitem}
\usepackage{fancyhdr}
\usepackage{xcolor}

% ===================== HEADER & FOOTER =====================
\pagestyle{fancy}
\fancyhf{}
\lhead{CSE 400: Fundamentals of Probability in Computing}
\rhead{Lecture 5 Scribe}
\cfoot{\thepage}

% ===================== TITLE =====================
\title{
\normalsize School of Engineering and Applied Science (SEAS), Ahmedabad University\\
\vspace{0.2cm}
\textbf{CSE 400: Fundamentals of Probability in Computing}\\
\Large Lecture 5: Bayes' Theorem, Random Variables, and PMF
}
\author{}

\date{January 20, 2026}

\begin{document}
\maketitle
\vspace{-1.5cm}

\begin{center}
\begin{tabular}{ll}
\textbf{Instructor:} & Dhaval Patel, PhD \\[1ex]
\textbf{Purpose:} & Exam-oriented reference material
\end{tabular}
\end{center}

\hrule
\vspace{0.5cm}

% ===================== CONTENT =====================

\section{Bayes' Theorem}

\subsection{Weighted Average of Conditional Probabilities}
Let $A$ and $B$ be events. Event $A$ can be decomposed as
\[
A = (A \cap B) \cup (A \cap B^c),
\]
where the events $(A \cap B)$ and $(A \cap B^c)$ are mutually exclusive.  
By Axiom 3 (additivity),
\[
\Pr(A) = \Pr(A \cap B) + \Pr(A \cap B^c).
\]

Using conditional probability,
\[
\Pr(A \cap B) = \Pr(A \mid B)\Pr(B), \quad
\Pr(A \cap B^c) = \Pr(A \mid B^c)\Pr(B^c).
\]

Hence,
\[
\Pr(A) = \Pr(A \mid B)\Pr(B) + \Pr(A \mid B^c)[1 - \Pr(B)].
\]
Thus, $\Pr(A)$ is a weighted average of conditional probabilities.

\subsection{Example 3.1 (Part 1)}
An insurance company classifies people as accident-prone ($A$) or not ($A^c$).  
Let $B$ denote the event that a policyholder has an accident within one year.  
Using total probability,
\[
\Pr(B) = \Pr(B \mid A)\Pr(A) + \Pr(B \mid A^c)\Pr(A^c).
\]

\subsection{Example 3.1 (Part 2)}
Given that an accident occurred, we compute
\[
\Pr(A \mid B) = \frac{\Pr(B \mid A)\Pr(A)}{\Pr(B)}.
\]

\section{Law of Total Probability and Bayes Formula}

Let $\{B_1, B_2, \dots, B_n\}$ be a partition of the sample space with $\Pr(B_i) > 0$.

\subsection*{Law of Total Probability}
\[
\Pr(A) = \sum_{i=1}^{n} \Pr(A \mid B_i)\Pr(B_i).
\]

\subsection*{Bayes Formula (Proposition 3.1)}
\[
\Pr(B_i \mid A) =
\frac{\Pr(A \mid B_i)\Pr(B_i)}
{\sum_{j=1}^{n} \Pr(A \mid B_j)\Pr(B_j)}.
\]

Here, $\Pr(B_i)$ is the \textit{a priori} probability and $\Pr(B_i \mid A)$ is the \textit{posteriori} probability.

\subsection{Example 3.2 (Cards Problem)}
Three cards: RR, RB, BB, each chosen with probability $1/3$.  
Let $R$ denote the event that the upper side is red.
\[
\Pr(R) = \Pr(R \mid RR)\Pr(RR) + \Pr(R \mid RB)\Pr(RB) + \Pr(R \mid BB)\Pr(BB).
\]
\[
\Pr(R) = (1)\left(\frac{1}{3}\right) + \left(\frac{1}{2}\right)\left(\frac{1}{3}\right) + (0)\left(\frac{1}{3}\right)
= \frac{1}{2}.
\]

\section{Random Variables}
A random variable is a real-valued function defined on the sample space.  
Interest often lies in a function of the outcome rather than the outcome itself.

\subsection*{Examples}
\begin{itemize}
\item Sum of two dice
\item Number of heads in repeated coin tosses
\end{itemize}

The distribution can be visualized using a bar diagram with height $\Pr(X = x)$ at value $x$.

\subsection{Example: Three Coin Tosses}
Let $X$ denote the number of heads in three fair coin tosses.

\section{Probability Mass Function (PMF)}
A random variable taking at most countably many values is called discrete.

Let $X$ have range $R_X = \{x_1, x_2, \dots\}$.  
The probability mass function is defined as
\[
p(x_k) = \Pr(X = x_k).
\]

\subsection*{Properties}
\[
p(x_k) \ge 0, \quad \sum_k p(x_k) = 1.
\]

\subsection{PMF Example}
Given a random variable $Y$ with
\[
\Pr(Y=0)=\frac{1}{8}, \quad \Pr(Y=1)=\frac{3}{8}, \quad
\Pr(Y=2)=\frac{3}{8}, \quad \Pr(Y=3)=\frac{1}{8},
\]
we verify
\[
\sum_{k=0}^{3} \Pr(Y=k) = 1.
\]

\subsection{Poisson Distribution}
Let
\[
p(i) = c\frac{\lambda^i}{i!}, \quad i=0,1,2,\dots, \ \lambda>0.
\]
Since
\[
\sum_{i=0}^{\infty} p(i) = 1 \implies c e^{\lambda} = 1 \implies c = e^{-\lambda},
\]
we obtain
\[
\Pr(X=0)=e^{-\lambda}, \quad
\Pr(X>2)=1-\sum_{i=0}^{2} e^{-\lambda}\frac{\lambda^i}{i!}.
\]

\vspace{0.5cm}
\hrule
\vspace{0.2cm}
\begin{center}
\small \textit{End of Lecture 5 Scribe}
\end{center}

\end{document}
